% GENERATED FILE. Edit canonical lessons, then run npm run generate:books.
% canonical-source-hash: fnv1a64:f51090dbb7b1c79d
% canonical-lessons: MR-C327-praja, MR-C327-rajmudra, MR-C327-nisan, MR-C327-sardar, MR-C327-vasavne, MR-R327-review-320-327

\chapter{Ruling the Land}
\label{ch:mr-ruling-the-land}
\hlchaptermodality{327}

\begin{chapteropening}
\textbf{By the end of this chapter:} \emph{I can talk about a ruler's subjects, seal, banner and new towns.}

Complete the last lesson of chapter 327: I can talk about a ruler's subjects, seal, banner and new towns.
\end{chapteropening}

\section[\texorpdfstring{\mr{प्रजा} (prajā)}{प्रजा (prajā)}]{\mr{प्रजा} (prajā) — subjects; the people of a king}

\label{lesson:MR-C327-praja}



\begin{quote}
Before the new one: say the Marathi for a princess, then the Marathi for a minister.
\end{quote}

\subsection*{You'll want to know: \mr{प्रजा}}
\textbf{\mr{प्रजा}} — \emph{prajā} — \textquotedblleft{}subjects; the people of a king\textquotedblright{}.

\begin{quote}
\mr{तो} \mr{राजा} \mr{प्रजेची} \mr{नेहमी} \mr{काळजी} \mr{घ्यायचा}.
\end{quote}

\begin{quote}
\emph{to rājā prajecī nehmī kāḷjī ghyāycā.}
\end{quote}

\begin{quote}
That king always took care of his subjects.
\end{quote}

Feminine.

Singular even for a whole people: \mr{प्रजा} \mr{सुखी} \mr{होती}.

\begin{cousinweb}
\textbf{\mr{प्रजा}} (\emph{prajā}) is a tatsama, taken from Sanskrit \emph{prajā}, 'offspring; people, subjects' (Monier-Williams).
\end{cousinweb}

\begin{grammarlens}[title={What you've built}]
One more word for telling what happened.
\end{grammarlens}

\subsection*{Your turn}
\begin{itemize}
\raggedright
  \item \emph{Say it:} \emph{prajā}
  \item \emph{Say it:} \emph{prajā}, once more
  \item \emph{Say it:} \emph{to rājā prajecī nehmī kāḷjī ghyāycā.} — That king always took care of his subjects.
  \item \emph{Recall:} say the Marathi for a princess, then the Marathi for a minister, then say \emph{prajā} again
\end{itemize}

\begin{tcolorbox}[breakable,colback=teal!4,colframe=teal!35!black,title={Before you move on}]
What does \mr{प्रजा} mean? (\textquotedblleft{}Subjects; the people of a king\textquotedblright{}.) Say it once more.
\end{tcolorbox}
\section[\texorpdfstring{\mr{राजमुद्रा} (rājmudrā)}{राजमुद्रा (rājmudrā)}]{\mr{राजमुद्रा} (rājmudrā) — a royal seal}

\label{lesson:MR-C327-rajmudra}



\begin{quote}
Before the new one: say the Marathi for a minister, then the Marathi for subjects.
\end{quote}

\subsection*{You'll want to know: \mr{राजमुद्रा}}
\textbf{\mr{राजमुद्रा}} — \emph{rājmudrā} — \textquotedblleft{}a royal seal\textquotedblright{}.

\begin{quote}
\mr{जुन्या} \mr{पत्रांवर} \mr{राजमुद्रा} \mr{आहे}.
\end{quote}

\begin{quote}
\emph{junyā patrāṁvar rājmudrā āhe.}
\end{quote}

\begin{quote}
There is a royal seal on the old letters.
\end{quote}

Feminine.

\mr{राज}-, royal, + \mr{मुद्रा}, a seal.

\begin{cousinweb}
\textbf{\mr{राजमुद्रा}} (\emph{rājmudrā}) is built from Sanskrit \emph{rāja-}, 'royal', and \emph{mudrā}, 'a seal, a stamp' (Monier-Williams).
\end{cousinweb}

\begin{grammarlens}[title={What you've built}]
One more word for telling what happened.
\end{grammarlens}

\subsection*{Your turn}
\begin{itemize}
\raggedright
  \item \emph{Say it:} \emph{rājmudrā}
  \item \emph{Say it:} \emph{rājmudrā}, once more
  \item \emph{Say it:} \emph{junyā patrāṁvar rājmudrā āhe.} — There is a royal seal on the old letters.
  \item \emph{Recall:} say the Marathi for a minister, then the Marathi for subjects, then say \emph{rājmudrā} again
\end{itemize}

\begin{tcolorbox}[breakable,colback=teal!4,colframe=teal!35!black,title={Before you move on}]
What does \mr{राजमुद्रा} mean? (\textquotedblleft{}A royal seal\textquotedblright{}.) Say it once more.
\end{tcolorbox}
\section[\texorpdfstring{\mr{निशाण} (niśāṇ)}{निशाण (niśāṇ)}]{\mr{निशाण} (niśāṇ) — a banner; a flag}

\label{lesson:MR-C327-nisan}



\begin{quote}
Before the new one: say the Marathi for subjects, then the Marathi for a royal seal.
\end{quote}

\subsection*{You'll want to know: \mr{निशाण}}
\textbf{\mr{निशाण}} — \emph{niśāṇ} — \textquotedblleft{}a banner; a flag\textquotedblright{}.

\begin{quote}
\mr{किल्ल्यावर} \mr{उंच} \mr{निशाण} \mr{फडकत} \mr{होतं}.
\end{quote}

\begin{quote}
\emph{killyāvar uñc niśāṇ phaḍkat hotaṁ.}
\end{quote}

\begin{quote}
A tall banner was flying over the fort.
\end{quote}

Neuter.

\mr{निशाण} \mr{फडकणे}, a flag flies.

\begin{cousinweb}
\textbf{\mr{निशाण}} (\emph{niśāṇ}) is a loan from Persian \emph{nišān}, 'a sign, a mark, a standard'.
\end{cousinweb}

\begin{grammarlens}[title={What you've built}]
One more word for telling what happened.
\end{grammarlens}

\subsection*{Your turn}
\begin{itemize}
\raggedright
  \item \emph{Say it:} \emph{niśāṇ}
  \item \emph{Say it:} \emph{niśāṇ}, once more
  \item \emph{Say it:} \emph{killyāvar uñc niśāṇ phaḍkat hotaṁ.} — A tall banner was flying over the fort.
  \item \emph{Recall:} say the Marathi for subjects, then the Marathi for a royal seal, then say \emph{niśāṇ} again
\end{itemize}

\begin{tcolorbox}[breakable,colback=teal!4,colframe=teal!35!black,title={Before you move on}]
What does \mr{निशाण} mean? (\textquotedblleft{}A banner; a flag\textquotedblright{}.) Say it once more.
\end{tcolorbox}
\section[\texorpdfstring{\mr{सरदार} (sardār)}{सरदार (sardār)}]{\mr{सरदार} (sardār) — a chieftain; a noble}

\label{lesson:MR-C327-sardar}



\begin{quote}
Before the new one: say the Marathi for a royal seal, then the Marathi for a banner.
\end{quote}

\subsection*{You'll want to know: \mr{सरदार}}
\textbf{\mr{सरदार}} — \emph{sardār} — \textquotedblleft{}a chieftain; a noble\textquotedblright{}.

\begin{quote}
\mr{राजाने} \mr{शूर} \mr{सरदाराला} \mr{एक} \mr{गाव} \mr{भेट} \mr{दिलं}.
\end{quote}

\begin{quote}
\emph{rājāne śūr sardārālā ek gāv bheṭ dilaṁ.}
\end{quote}

\begin{quote}
The king gave the brave chieftain a village as a gift.
\end{quote}

Masculine.

\mr{सरदाराला}, to the chieftain.

\begin{cousinweb}
\textbf{\mr{सरदार}} (\emph{sardār}) is a loan from Persian \emph{sardār}, 'a chief, a commander', built from \emph{sar}, 'head', and \emph{-dār}, 'one who holds'.
\end{cousinweb}

\begin{grammarlens}[title={What you've built}]
One more word for telling what happened.
\end{grammarlens}

\subsection*{Your turn}
\begin{itemize}
\raggedright
  \item \emph{Say it:} \emph{sardār}
  \item \emph{Say it:} \emph{sardār}, once more
  \item \emph{Say it:} \emph{rājāne śūr sardārālā ek gāv bheṭ dilaṁ.} — The king gave the brave chieftain a village as a gift.
  \item \emph{Recall:} say the Marathi for a royal seal, then the Marathi for a banner, then say \emph{sardār} again
\end{itemize}

\begin{tcolorbox}[breakable,colback=teal!4,colframe=teal!35!black,title={Before you move on}]
What does \mr{सरदार} mean? (\textquotedblleft{}A chieftain; a noble\textquotedblright{}.) Say it once more.
\end{tcolorbox}
\section[\texorpdfstring{\mr{वसवणे} (vasavṇe)}{वसवणे (vasavṇe)}]{\mr{वसवणे} (vasavṇe) — to found a settlement}

\label{lesson:MR-C327-vasavne}



\begin{quote}
Before the new one: say the Marathi for a banner, then the Marathi for a chieftain.
\end{quote}

\subsection*{You'll want to know: \mr{वसवणे}}
\textbf{\mr{वसवणे}} — \emph{vasavṇe} — \textquotedblleft{}to found a settlement\textquotedblright{}.

\begin{quote}
\mr{नदीकाठी} \mr{राजाने} \mr{एक} \mr{नवं} \mr{शहर} \mr{वसवलं}.
\end{quote}

\begin{quote}
\emph{nadīkāṭhī rājāne ek navaṁ śahar vasavlaṁ.}
\end{quote}

\begin{quote}
The king founded a new city on the riverbank.
\end{quote}

Past, with -\mr{ने} (he …): \mr{त्याने} \mr{वसवलं}, \emph{tyāne vasavlaṁ}.

Transitive: \mr{शहर} \mr{वसवणे}, to found a town; the place itself \mr{वसतो}, settles.

\begin{cousinweb}
\textbf{\mr{वसवणे}} (\emph{vasavṇe}) is a Marathi causative, 'to make settle', built from the inherited stem \emph{vas-} with \emph{-av-}: CDIAL lists Marathi \emph{vasṇẽ}, 'settles', under Sanskrit \emph{vasati}, 'stays, dwells' (CDIAL 11435).
\end{cousinweb}

\begin{grammarlens}[title={What you've built}]
One more word for telling what happened.
\end{grammarlens}

\subsection*{Your turn}
\begin{itemize}
\raggedright
  \item \emph{Say it:} \emph{vasavṇe}
  \item \emph{Say it:} \emph{vasavṇe}, once more
  \item \emph{Say it:} \emph{nadīkāṭhī rājāne ek navaṁ śahar vasavlaṁ.} — The king founded a new city on the riverbank.
  \item \emph{Recall:} say the Marathi for a banner, then the Marathi for a chieftain, then say \emph{vasavṇe} again
\end{itemize}

\begin{tcolorbox}[breakable,colback=teal!4,colframe=teal!35!black,title={Before you move on}]
What does \mr{वसवणे} mean? (\textquotedblleft{}To found a settlement\textquotedblright{}.) Say it once more.
\end{tcolorbox}
\section[(dialogue)]{Looking back: from Festival Days to Ruling the Land}

\label{lesson:MR-R327-review-320-327}



\begin{quote}
Say the Marathi for to decorate and a festivity.
\end{quote}

\subsection*{Your turn}
Say these aloud:

\begin{itemize}
\raggedright
  \item to decorate, a festivity, a grand ceremony, a pilgrimage, an auspicious time
  \item a row of diners, hosting guests, a get-together, to attend, attendance
  \item an award, to felicitate, an honour, to congratulate, to inaugurate
  \item a farewell, retirement, a housewarming, an anniversary, a closing
  \item a passing, a tribute, to organise, to go off, to compose
  \item a king, a queen, a palace, a royal court, to reign
  \item an emperor, a throne, a crown, a princess, a minister
  \item subjects, a royal seal, a banner, a chieftain, to found a settlement
\end{itemize}

\begin{tcolorbox}[breakable,colback=teal!4,colframe=teal!35!black,title={Before you move on}]
To decorate? (\textbf{\mr{सजवणे}}, \emph{sajavṇe}.) A festivity? (\textbf{\mr{उत्सव}}, \emph{utsav}.) A grand ceremony? (\textbf{\mr{सोहळा}}, \emph{sohḷā}.) A pilgrimage? (\textbf{\mr{यात्रा}}, \emph{yātrā}.) An auspicious time? (\textbf{\mr{मुहूर्त}}, \emph{muhūrta}.) A row of diners? (\textbf{\mr{पंगत}}, \emph{paṅgat}.) Hosting guests? (\textbf{\mr{पाहुणचार}}, \emph{pāhuṇcār}.) A get-together? (\textbf{\mr{स्नेहसंमेलन}}, \emph{snehsammelan}.) To attend? (\textbf{\mr{उपस्थित} \mr{राहणे}}, \emph{upasthit rāhṇe}.) Attendance? (\textbf{\mr{हजेरी}}, \emph{hajerī}.) An award? (\textbf{\mr{पुरस्कार}}, \emph{puraskār}.) To felicitate? (\textbf{\mr{सत्कार} \mr{करणे}}, \emph{satkār karṇe}.) An honour? (\textbf{\mr{गौरव}}, \emph{gaurav}.) To congratulate? (\textbf{\mr{अभिनंदन} \mr{करणे}}, \emph{abhinandan karṇe}.) To inaugurate? (\textbf{\mr{उद्घाटन} \mr{करणे}}, \emph{udghāṭan karṇe}.) A farewell? (\textbf{\mr{निरोप}}, \emph{nirop}.) Retirement? (\textbf{\mr{निवृत्ती}}, \emph{nivṛttī}.) A housewarming? (\textbf{\mr{गृहप्रवेश}}, \emph{gṛhapraveś}.) An anniversary? (\textbf{\mr{वर्धापनदिन}}, \emph{vardhāpandin}.) A closing? (\textbf{\mr{समारोप}}, \emph{samārop}.) A passing? (\textbf{\mr{निधन}}, \emph{nidhan}.) A tribute? (\textbf{\mr{श्रद्धांजली}}, \emph{śraddhāñjalī}.) To organise? (\textbf{\mr{आयोजन} \mr{करणे}}, \emph{āyojan karṇe}.) To go off? (\textbf{\mr{पार} \mr{पडणे}}, \emph{pār paḍṇe}.) To compose? (\textbf{\mr{रचणे}}, \emph{racṇe}.) A king? (\textbf{\mr{राजा}}, \emph{rājā}.) A queen? (\textbf{\mr{राणी}}, \emph{rāṇī}.) A palace? (\textbf{\mr{राजवाडा}}, \emph{rājvāḍā}.) A royal court? (\textbf{\mr{दरबार}}, \emph{darbār}.) To reign? (\textbf{\mr{राज्य} \mr{करणे}}, \emph{rājya karṇe}.) An emperor? (\textbf{\mr{सम्राट}}, \emph{samrāṭ}.) A throne? (\textbf{\mr{सिंहासन}}, \emph{siṁhāsan}.) A crown? (\textbf{\mr{मुकुट}}, \emph{mukuṭ}.) A princess? (\textbf{\mr{राजकन्या}}, \emph{rājkanyā}.) A minister? (\textbf{\mr{प्रधान}}, \emph{pradhān}.) Subjects? (\textbf{\mr{प्रजा}}, \emph{prajā}.) A royal seal? (\textbf{\mr{राजमुद्रा}}, \emph{rājmudrā}.) A banner? (\textbf{\mr{निशाण}}, \emph{niśāṇ}.) A chieftain? (\textbf{\mr{सरदार}}, \emph{sardār}.) To found a settlement? (\textbf{\mr{वसवणे}}, \emph{vasavṇe}.)
\end{tcolorbox}
